\section{Ensembles and CRPS weighting --- why averaging models wins}

\subsection{The idea}

An ensemble averages the forecasts of several models.
It wins when the members make \emph{different} mistakes.
If model A is 2 EUR too high and model B is 2 EUR too low,
their average is perfect. Correlated errors average to nothing.
So the recipe is: combine models that are individually good
AND structurally different.

Our members: LGBM (trained, covariates), LEAR (linear, covariates),
Chronos (zero-shot, univariate). Three different failure modes.

\subsection{CRPS and our crps3 proxy}

CRPS (continuous ranked probability score) measures a whole
predictive distribution against one observed value.
Lower is better. It generalizes MAE to distributions.

We store only three quantiles (P10/P50/P90), not a full distribution.
So we use the mean pinball loss over those three quantiles.
We call it \texttt{crps3} --- the name admits it is a proxy.

\subsection{Inverse-score weights, past-only}

Weight for model $m$ on day $D$:
\[
w_m(D) \propto \frac{1}{\overline{\text{crps3}}_m(D-60 \dots D-1)}
\]
Better recent skill $\rightarrow$ bigger weight. The window sees only
days before $D$ --- no leakage. The first 60 days use equal weights
(no history yet). Weights adapt when a model degrades.

\subsection{Quantile-wise blending}

We average each quantile separately: blended P10 is the weighted
mean of member P10s, and so on. Then we re-sort the three numbers
per hour so P10 $\le$ P50 $\le$ P90 always holds.
Caveat: averaging quantiles is not the same as averaging
distributions (that would be Vincentization vs mixture). For three
quantiles the practical difference is small.

\subsection{The double conformal trick}

Averaging bands from three calibrated models produced an
over-covering band: 84.2\% instead of the nominal 80\%.
Why: member errors are positively but not perfectly correlated,
so the blended band is wider than needed.
Fix: run rolling CQR \emph{again} on the blended band.
The conformal offset $Q$ goes negative and tightens the band.
Result: coverage 79.9\%, Winkler improved, MAE untouched.

\subsection{Worked example: our numbers}

\begin{tabular}{lccc}
\hline
Model & MAE & Coverage & Winkler \\
\hline
Ensemble + CQR & 17.34 & 79.9\% & 84.7 \\
Ensemble raw & 17.34 & 84.2\% & 85.2 \\
Equal weights & 17.46 & 84.3\% & 85.9 \\
Champion LGBM & 17.87 & 78.9\% & 89.5 \\
\hline
\end{tabular}

\vspace{0.5em}
Note the equal-weight row. Skill weighting adds only 0.12 MAE over
equal weights. Most of the gain is the diversity itself.
This is a standard finding in the combination literature.

\subsection{Interview line}

``I blended LGBM, LEAR and a zero-shot Chronos with inverse-CRPS
weights from a trailing 60-day window. The blend beat the best
single model by 0.53 MAE, DM $p < 0.001$. The surprise: Chronos is
4 MAE worse alone, but it still helps --- its errors are
uncorrelated with the structural models. Diversity, not weighting,
does the work.''
